\section{Fundamentação teórica e trabalhos relacionados}\label{sec:fundamentacao}

\subsection{O Model Context Protocol}

O MCP é um protocolo aberto que padroniza a comunicação entre aplicações de IA
e fontes externas de contexto e execução \citep{anthropic2024mcp}. Sua
especificação define uma arquitetura cliente--servidor na qual o
\emph{host} (por exemplo, um editor de código ou assistente conversacional)
instancia um cliente MCP por servidor conectado. A comunicação emprega
JSON-RPC 2.0 e pode ocorrer por entrada/saída padrão (\emph{stdio}) ou por
transportes HTTP com streaming \citep{mcp2025spec}.

O protocolo define três primitivas principais que um servidor pode oferecer
\citep{mcp2025docs}:

\begin{itemize}[leftmargin=1.4em,itemsep=0.2em]
  \item \textbf{Tools}: funções executáveis, descritas por esquema JSON, que o
        modelo pode invocar para produzir efeitos (criar arquivos, modificar
        geometria, executar verificações);
  \item \textbf{Resources}: dados somente leitura identificados por URI, como o
        estado de um projeto ou listas de componentes;
  \item \textbf{Prompts}: modelos de interação predefinidos que orientam o
        uso do servidor para tarefas recorrentes.
\end{itemize}

A especificação utilizada nesta implementação é a de 18 de junho de 2025
\citep{mcp2025spec}, que consolida o ciclo de vida de inicialização,
negociação de capacidades e transporte. A adoção do MCP por diferentes
fornecedores de modelos e editores reduziu o acoplamento entre assistentes e
ferramentas, tornando viável publicar um único servidor para todo um domínio,
como o projeto de PCBs.

\subsection{Agentes baseados em LLMs e uso de ferramentas}

A capacidade de um agente de interagir com o ambiente depende não apenas do
modelo, mas do desenho da interface que lhe é oferecida. \citet{yao2023react}
demonstram que intercalar raciocínio e ação melhora o desempenho em tarefas
que exigem decisões encadeadas. \citet{schick2023toolformer} e
\citet{patil2023gorilla} evidenciam que modelos podem aprender a invocar APIs
e ferramentas externas com acurácia crescente quando as chamadas são expostas
de forma estruturada. \citet{yang2024sweagent} mostram, no contexto de
engenharia de software, que a interface agente--computador (ACI) tem impacto
direto no desempenho: nomes de comandos claros, saídas concisas e
descobribilidade de ferramentas importam tanto quanto o modelo utilizado.

Esses resultados motivam decisões de projeto adotadas neste trabalho: exposição
individual de cada ferramenta (sem roteador oculto), descrições objetivas,
mecanismos explícitos de descoberta e respostas padronizadas e testáveis.

\subsection{Automação de EDA e o KiCad}

O KiCad é uma suíte de EDA de código aberto que abrange captura esquemática,
layout de placa, bibliotecas de símbolos e \emph{footprints}, visualização 3D e
geração de arquivos de fabricação \citep{kicaddocs}. Seus arquivos de projeto
são expressões-S (\emph{S-expressions}) versionáveis em texto, o que facilita
automação e revisão por ferramentas externas.

A extensibilidade do KiCad se apoia em três mecanismos complementares:
(i) a API Python \texttt{pcbnew}, disponibilizada por meio de \emph{bindings}
SWIG \citep{swig}; (ii) a API IPC, que permite comunicação em tempo real com a
interface gráfica; e (iii) o utilitário de linha de comando \texttt{kicad-cli},
que executa exportações e verificações de forma independente da interface
gráfica \citep{kicadcli,kicaddev}. Esses mecanismos são a base de integrações
automatizadas, inclusive as conduzidas por agentes de IA.

\subsection{Trabalhos relacionados}

A aplicação de LLMs ao domínio de \emph{design} de circuitos tem exemplos
relevantes. \citet{liu2023chipnemo} adaptam modelos ao domínio de projeto de
\emph{chips} industriais (assistente de engenharia, geração de \emph{scripts}
de EDA e análise de defeitos), demonstrando ganhos com especialização de
domínio. Trabalhos dessa natureza, no entanto, concentram-se na adaptação do
\emph{modelo} a um domínio; a integração entre assistentes genéricos e
ferramentas de EDA permanece fragmentada, tipicamente resolvida por
\emph{scripts} específicos e não reutilizáveis.

Do lado das ferramentas, a automação do KiCad por \emph{scripts} Python é
consolidada, porém cada integração reproduz decisões de sessão, tratamento de
erros e descoberta de operações. O MCP altera esse cenário ao oferecer um
contrato único entre assistentes e ferramentas \citep{mcp2025spec}. Neste
contexto, o \texttt{mcp-kicad-vscode} se posiciona como um servidor MCP de
propósito geral para o fluxo de PCB no KiCad, com cobertura ampla,
integrações de fabricação e suíte de testes própria --- combinação ainda
escassa na literatura e no ecossistema de código aberto.

A lacuna endereçada é, portanto, dupla: (i) ausência de uma camada padronizada
e descobrível que cubra o ciclo de projeto de PCB de ponta a ponta
(esquemático, layout, verificação, fabricação); e (ii) ausência de evidências
publicadas sobre robustez e manutenibilidade desse tipo de integração,
incluindo o comportamento sob operações destrutivas e a cobertura de testes.
